% ---------------------------------------------------------------------
%  04 Method
%  Mirrors the repo pipeline:
%    src/tasks/ball_detection, src/tasks/court_detection,
%    src/submodules (2D pose + GVHMR), src/tasks/blcs, src/tasks/plcs
% ---------------------------------------------------------------------
\section{Method}
\label{sec:method}

\subsection{Overview and problem statement}
\label{subsec:overview}

Given $\numcams$ synchronised video streams of a tennis rally, each of
$\numframes$ frames, we seek the 3D ball trajectory
$\{\ballThreeD{\timestep}\}_{\timestep=1}^{\numframes}$ and, for each player,
a root transform $(\rootrot{\timestep}, \rootpos{\timestep}) \in \SEthree$
together with articulated body pose
$\poseThreeD{\timestep}$, all expressed in a common court coordinate system.

\Todo{State the coordinate convention explicitly (origin, axes, units) and
point to \cref{fig:pipeline}. Define what is assumed known (court dimensions)
versus estimated (camera pose via court homography).}

% --- Pipeline figure --------------------------------------------------
\begin{figure}[t]
  \centering
  % \includegraphics[width=\linewidth]{assets/figures/pipeline.pdf}
  \fbox{\parbox[c][0.22\textheight][c]{\linewidth}{\centering
    \Todo{pipeline.pdf --- 2D perception (ball / court / pose) $\rightarrow$
    lifting (\BLCS{}, \PLCS{}) $\rightarrow$ fusion with \GVHMR{}}}}
  \caption{\Todo{Pipeline caption. Name every box and every arrow's tensor
  shape; the reader should be able to reimplement the interfaces from this
  figure alone.}}
  \label{fig:pipeline}
\end{figure}

\subsection{Stage 1: per-view 2D perception}
\label{subsec:perception}

\subsubsection{Ball detection}
\label{subsubsec:ball_detection}
\Todo{Backbone, heatmap head, temporal context ($T$ frames in, 1 heatmap out),
sub-pixel refinement, and the visibility gate. Give the loss.}

\begin{equation}
  \label{eq:ball_loss}
  \mathcal{L}_{\mathrm{ball}} = \Todo{heatmap loss + subpixel offset + visibility}
\end{equation}

\subsubsection{Court keypoint detection}
\label{subsubsec:court_detection}
\Todo{$\numcourtkp = 14$ ground-plane keypoints; per-keypoint visibility;
recovery of the homography $\homog{\cam}$ by robust fitting, and how
$\homog{\cam}$ yields the camera projection $\proj{\cam}$.}

\subsubsection{Player detection and 2D pose}
\label{subsubsec:pose_2d}
\Todo{Detector + top-down 2D pose over $\numjoints$ joints; track association
across frames and across views.}

\subsection{Stage 2: lifting to court coordinates}
\label{subsec:lifting}

\subsubsection{\BLCS{}: ball lifting}
\label{subsubsec:blcs}
\Todo{Inputs: per-view 2D ball detections $\ballTwoD{\cam,\timestep}$ and court
geometry. Architecture. Why a learned lifter rather than pure triangulation
(missed detections, single-view intervals, motion blur). Physics prior /
smoothness. Loss.}

\begin{equation}
  \label{eq:blcs}
  \ballThreeD{1:\numframes} = f_{\BLCS{}}\!\left(
    \{\ballTwoD{\cam,1:\numframes}\}_{\cam=1}^{\numcams},\;
    \{\homog{\cam}\}_{\cam=1}^{\numcams}
  \right)
\end{equation}

\subsubsection{\PLCS{}: player root position and orientation}
\label{subsubsec:plcs}
\Todo{Inputs: 2D pose sequences and court keypoints. Output: root translation
$\rootpos{\timestep}$ and rotation $\rootrot{\timestep}$. Rotation
parameterisation and the angle loss that removes the $180^\circ$ flip failure
mode. Separate trunks / auxiliary heads --- justify each with the ablation it
is supported by.}

\begin{equation}
  \label{eq:plcs}
  (\rootrot{1:\numframes}, \rootpos{1:\numframes}) = f_{\PLCS{}}\!\left(
    \{\poseTwoD{\cam,1:\numframes}\}_{\cam=1}^{\numcams},\;
    \{\courtkp{\cam}\}_{\cam=1}^{\numcams}
  \right)
\end{equation}

\subsection{Stage 3: fusion into a court-grounded reconstruction}
\label{subsec:fusion}

\Todo{\GVHMR{} supplies body pose in its own gravity-aligned frame; \PLCS{}
supplies the court-frame root transform. Give the composition explicitly,
including how the two frames are reconciled and what happens when they
disagree. This is where the paper's central claim lives --- do not hand-wave
it.}

\subsection{Training data}
\label{subsec:training_data}
\Todo{Simulated rallies (generation procedure, randomisation) and real
annotated clips. State the sim-to-real design decision: \PLCS{} is trained with
$\numcourtkp = 14$ so that its input distribution matches the real court
detector's output.}

\subsection{Implementation details}
\label{subsec:implementation}
\Todo{Backbones, optimiser, schedule, batch size, hardware, wall-clock. Keep
this short; the full configuration goes in \cref{sec:appendix_config}.}
